% Resumen de docs/adr/ADR-002-mecanismo-comunicacion.md. Responsable: Mario (Líder de Desarrollo).
\subsection{ADR-002: mecanismo de comunicación principal}

\textbf{Decisión.} gRPC para las llamadas directas entre servicios. Hoy hay una sola: antes de abrir la rúbrica, Evaluación y actas le pregunta a Solicitudes y documentos si el informe y el reporte de URKUND están, y cuál es el porcentaje. El contrato está en \nolinkurl{src/proto/solicitudes.proto} y lo comparten los dos servicios. REST queda para la aplicación web y el API Gateway, y los eventos van por RabbitMQ (sección~\ref{sec:comunicacion}).

\textbf{Por qué.} Comparamos las dos opciones que medimos en PE-U1, sockets TCP y gRPC, en una sola máquina y sin Docker. En la corrida oficial del 22 de noviembre de 2026, con un cliente y mensajes de 100~bytes, la mediana de ida y vuelta fue de 0,045~ms con TCP y de 0,291~ms con gRPC, unas 6,5 veces más. Con 10~KB la brecha bajó a 4,2 veces, y con tres clientes a la vez gRPC duplicó su rendimiento (de 2640 a 5636 mensajes por segundo), aunque TCP siguió adelante en números absolutos. O sea, TCP es más rápido, pero la diferencia es de un cuarto de milisegundo en una consulta que se hace una vez antes de abrir la rúbrica, con 30 usuarios a la vez como máximo (RNF-05). Eso no lo nota nadie. Lo que sí importa es que el dato que decide si el tribunal califica llegue completo y con el tipo correcto, y con gRPC eso lo revisa el código generado desde el \texttt{.proto}. Con sockets, en cambio, tuvimos que validar a mano cada campo del JSON. Además, el equipo ya había usado gRPC en PE-U1.

\textbf{Opción descartada.} Sockets TCP con mensajes JSON por línea, porque el formato, los tiempos de espera y los errores había que programarlos a mano, y un cambio de campo en un servicio solo se notaba al ejecutar.

\textbf{Consecuencias.} Si Solicitudes no responde en 5~segundos, Evaluación recibe \texttt{UNAVAILABLE} o \texttt{DEADLINE\_EXCEEDED} y no abre la rúbrica, como pide la sección~\ref{sec:cap-pacelc}. A cambio, hay que generar el código cada vez que cambia el \texttt{.proto}, la llamada no se puede probar desde el navegador y, en E1, el canal va sin cifrar porque todo corre en la misma máquina.

El registro completo está en \nolinkurl{docs/adr/ADR-002-mecanismo-comunicacion.md}.
